\documentclass[oneside,12pt]{exam}          % ----- exam

%%%%%%%%%%%%%%%%%%%%%%%%%%%%
%%%%% Police et langue %%%%%
%%%%%%%%%%%%%%%%%%%%%%%%%%%%

\usepackage[x11names,table]{xcolor}                           % ----- Couleurs
\usepackage[utf8]{inputenc}                             % ----- UTF-8
\usepackage[french]{babel}                              % ----- French
\usepackage{amsmath,amssymb,mathrsfs}               % ----- Maths
\usepackage{fontawesome5,txfonts}						% ----- Caractères / Symboles
\usepackage{eurosym}                                    % ----- Currencies

\usepackage{helvet}										% ----- Police texte
\renewcommand{\familydefault}{\sfdefault}
\usepackage[T1]{fontenc}

\usepackage{mathpazo}									% ----- Police maths


%%%%%%%%%%%%%%%%
%%%%% TikZ %%%%%
%%%%%%%%%%%%%%%%

\usepackage{pgfplots}                                   % ----- Plots
    \pgfplotsset{compat=1.18}
\usepackage{tikz, tikz-3dplot, tkz-tab , tkz-euclide}                 % ----- TikZ, TikZ 3D, Tableaux

    \usetikzlibrary{math,arrows,arrows.meta}                            % ----- Maths, Vecteurs
    \usetikzlibrary{automata, positioning,patterns.meta}    % ----- Graphs
    \usetikzlibrary{decorations.fractals}                   % ----- Fractals
    \usetikzlibrary{tikzmark}
    \usetikzlibrary{fadings}
    \usetikzlibrary{calc}
    \usetikzlibrary{angles}

\tikzset{arrow style/.style = {\tkzTabDefaultWritingColor,-Stealth,> = \tkzTabDefaultArrowStyle,shorten > = \tkzTabDefaultSep,shorten < = \tkzTabDefaultSep}}		% ----- Modifier flèche tabVar

\tikzset{node style/.style = {inner sep = \tkzTabDefaultSep, outer sep = \tkzTabDefaultSep}}

%%%%%%%%%%%%%%%%%%%%
%%%%% Pratique %%%%%
%%%%%%%%%%%%%%%%%%%%

\usepackage{titlesec}							% ----- Titres et sections
\usepackage{calc}                                       % ----- Calc / Minipage
%\usepackage{fancyhdr}                                   % ----- Header and footer
\usepackage{geometry}                                   % ----- Géométrie de la page
\usepackage{tcolorbox}                                  % ----- Boxes environment
    \tcbuselibrary{skins,listings,listingsutf8}             % ----- All libraries
\usepackage{varwidth}                                   % ----- Adjust box width
\usepackage{fourier-orns}                               % ----- Déco
\usepackage{witharrows}									% ----- Équations avec flèches
\usepackage{array}										% ----- Environnements array / tabular
\usepackage{enumitem}									% ----- Customisation des \enumerate
\usepackage[c]{esvect}									% ----- Vecteurs
\usepackage{diagbox}									% ----- Découper des cases de \tabular
\usepackage{esvect}
\usepackage{graphicx}									% ----- Insérer des images			% ----- Packages utiles
%%%%%%%%%%%%%%%%%%%%%%%%
%%%%% Mise en page %%%%%
%%%%%%%%%%%%%%%%%%%%%%%%

\usepackage{parskip}									% ----- Parskip
    \setlength{\parskip}{\medskipamount}
    \newcommand{\restoreskip}{\setlength{\parskip}{\medskipamount}}
    
\setlength\parindent{0mm}                               % ----- No indent

\usepackage{setspace}                                   % ----- Interligne
	\setstretch{0.95}                           % Interligne 0.95
\usepackage{makecell}



%%%%%%%%%%%%%%%%%%%%%%%%%%%%%%%
%%%%% Format des sections %%%%%
%%%%%%%%%%%%%%%%%%%%%%%%%%%%%%%

\titlespacing{\section}{0pt}{5mm}{0pt}
\titlespacing{\subsection}{0pt}{2mm}{1mm}


\titleformat{\section}
    {\large\normalfont\bfseries}
    {Exercice \no \thesection\, - \,}
    {0pt}
    {\large\bfseries}

\renewcommand{\thesubsection}{\Alph{subsection}}
\titleformat{\subsection}
    {\normalfont\bfseries}
    {Partie \thesubsection\,: \,}
    {0pt}
    {\bfseries}





		% ----- Préambule feuilles de TD
%%%%%%%%%%%%%%%%%%%%%%%%%%%%%
%%%% Lettres double barre %%%
%%%%%%%%%%%%%%%%%%%%%%%%%%%%%

\newcommand\CC{%
        {\mathbb C}} % ----- Ensemble C

\newcommand\RR{%
        {\mathbb R}} % ----- Ensemble R

\newcommand\QQ{%
        {\mathbb Q}} % ----- Ensemble Q

\newcommand\ZZ{%
        {\mathbb Z}} % ----- Ensemble Z

\newcommand\NN{%
        {\mathbb N}} % ----- Ensemble N

\newcommand\DD{%
        {\mathbb D}} % ----- Ensemble D
        
\newcommand\UU{%
        {\mathbb U}} % ----- Ensemble U

\newcommand\PP{%
        {\mathbb P}} % ----- P probabiliste

\newcommand\EE{%
        {\mathbb E}} % ----- E espérance

\newcommand\VV{%
        {\mathbb V}} % ----- V variance



%%%%%%%%%%%%%%%%%%%%%%%%%%%%%%%%
%%%% Lettres calligraphiques %%%
%%%%%%%%%%%%%%%%%%%%%%%%%%%%%%%%

\newcommand\Ccal{%
        {\mathscr{C}}} % ----- C calligraphique

\newcommand\Dcal{%
        {\mathscr{D}}} % ----- D calligraphique

\newcommand\Pcal{%
        {\mathscr{P}}} % ----- P calligraphique

\newcommand\Rcal{%
        {\mathscr{R}}} % ----- R calligraphique

\newcommand\Fcal{%
        {\mathscr{F}}} % ----- F calligraphique

\newcommand\Scal{%
        {\mathscr{S}}} % ----- S calligraphique

\newcommand\Acal{%
        {\mathscr{A}}} % ----- A calligraphique

\newcommand\Vcal{%
        {\mathscr{V}}} % ----- V calligraphique
        
\newcommand\Tcal{%
        {\mathscr{T}}} % ----- V calligraphique



%%%%%%%%%%%%%%%%%%%%%%%%%%%%%%%%%%%%%
%%% Symboles et/ou calculs usuels %%%
%%%%%%%%%%%%%%%%%%%%%%%%%%%%%%%%%%%%%

%%%%% Analyse %%%%%

\newcommand\limx{%
        {\lim\limits_{x\to +\infty}}\,\,} % ----- Limite x en + infini

\newcommand\limn{%
        {\lim\limits_{n\to +\infty}}\,\,} % ----- Limite n en + infini

\newcommand\limite[1]{%
        {\lim\limits_{\substack{#1}}}\,\,} % ----- Limite X vers Y

\newcommand\nN{%
        {n\in\mathbb N}} % ----- n dans N

\newcommand\xR{%
        {x\in\mathbb R}} % ----- x dans R

\newcommand\zC{%
        {z\in\mathbb C}} % ----- z dans C

\newcommand\somme[2]{%
        {\displaystyle\sum_{#1}^{#2}}}    % ----- Somme

\newcommand\sommek[2]{%
        {\displaystyle\sum_{k=#1}^{#2}}}    % ----- Somme k

\newcommand\inte[2]{%
        {\displaystyle\int_{#1}^{#2}}}   % ----- Intégrale

\newcommand\dd{%
        {\,\,\mathrm{d}}}       % ----- d différentiel

\newcommand\floor[1]{%
        {\lfloor #1 \rfloor}}   % ----- Partie entière inférieure

\newcommand\ceil[1]{%
        {\lceil #1 \rceil}}     % ----- Partie entière supérieure

\newcommand\abs[1]{%
        {\left\lvert #1 \right\rvert}}     % ----- Valeur absolue

\newcommand\pgcd{%
        {\mathrm{pgcd}}}       % ----- pgcd

\newcommand\ppcm{%
        {\mathrm{ppcm}}}       % ----- ppcm



%%%%% Probabilités %%%%%

\newcommand\barre[1]{%
    {\overline{#1}}} % ----- A barre
    
\newcommand\bbigcup{%
    \,\,\bigcup\,\,
    } % ----- Union avec espace 

\newcommand\bbigcap{%
    \,\,\bigcap\,\,
    } % ----- Intersection avec espace


%%%%% Géométrie %%%%%

\newcommand\vcoord[2]{			% ----- Coordonnées d'un vecteur (Nom + coord)
	\renewcommand{\arraystretch}{1.1} \arraycolsep=0.5\arraycolsep
	\vv{#1}\mathbin{=}\begin{pmatrix}#2\end{pmatrix}
} 

\newcommand\Oij{\left(O\,;\,\vv{i}\,;\,\vv{j}\right)} % ----- (O;i;j)

\def\Oijk{\left(O\,;\,\Vec{i}\,;\,\Vec{j}\,;\,\Vec{k}\right)} % ---- (O;i;j;k)

\newcommand\norme[1]{%
        {\left\lVert #1 \right\rVert}}     % ----- Norme

\newcommand\normev[1]{%
        {\left\lVert \vv{#1} \right\rVert}}     % ----- Norme d'un vecteur

\newcommand\scal[2]{%
    {\vv{#1}\cdot\vv{#2}}} % ----- Produit scalaire
    
\newcommand\smeq{{\,=\,}}	% ----- Symbole '=' avec petit espace



%%%%% Algèbre %%%%%

\newcommand\Reel{				% ----- Partie réelle
    {\mathfrak{Re}}} 

\newcommand\Imag{				% ----- Partie imaginaire
    {\mathfrak{Im}}} 

\newcommand\matrice[1]{{ 				% ----- Matrice
    \renewcommand{\arraystretch}{0.7}
    \arraycolsep=0.7\arraycolsep\ensuremath{\begin{pmatrix} #1 \end{pmatrix}}}}

\newcommand\Ouv{(O\,;\,\vv{u}\,;\,\vv{v})} 	% ----- Plan complexe (O;u;v)


                 % ----- Macros mathématiques
%%%%%%%%%%%%%%%%%%%%%%%%%%%%%%%%
%%%%% ----- Packages ----- %%%%%
%%%%%%%%%%%%%%%%%%%%%%%%%%%%%%%%

\usepackage{tcolorbox}                                  % ----- Boxes environment
    \tcbuselibrary{skins,listings,listingsutf8,breakable}             % ----- All libraries
\usepackage[bold=0.7]{xfakebold}


%%%%%%%%%%%%%%%%%%%%%%%%%%%%%%%%%%
%%%%% ----- New colors ----- %%%%%
%%%%%%%%%%%%%%%%%%%%%%%%%%%%%%%%%%

\definecolor{persianblue}{HTML}{221177} % ----- New blues
\definecolor{lightblue}{HTML}{EEFFFF}
\definecolor{softblue}{HTML}{6655BB}

\definecolor{maroon}{HTML}{992222}      % ----- New reds
\definecolor{lightred}{HTML}{FAEEEE}

\definecolor{calgreen}{HTML}{005500}    % ----- New greens
\definecolor{lightgreen}{HTML}{EEFFEE}
\definecolor{softgreen}{HTML}{449944}

\definecolor{brickorange}{HTML}{CC7033} % ----- New oranges
\definecolor{lightorange}{HTML}{FFF2E6}

\definecolor{borderpink}{HTML}{aa00bb} % ----- New pinks
\definecolor{lightpink}{HTML}{ffeeff}

\definecolor{newred}{HTML}{CC0000}     % ----- Colors for geometry
\definecolor{newblue}{HTML}{0000CC}
\definecolor{newgreen}{HTML}{00AA00}
\definecolor{neworange}{HTML}{FFAA00}







%%%%%%%%%%%%%%%%%%%%%%%%%%%%%%%%%%%%%%%%
%%%%% ----- \begin{theoreme} ----- %%%%%
%%%%%%%%%%%%%%%%%%%%%%%%%%%%%%%%%%%%%%%%

\newtcolorbox{theoreme}[2][]{%
    enhanced , breakable , parbox=false ,
    arc=5pt,
    frame style={left color=persianblue, right color=softblue},
    colback=lightblue,
    fonttitle=\bfseries,
    title={\faBook~~ \underline{Théorème} : #2~~\faBook},
    attach boxed title to top left={xshift=5mm,yshift=-2mm},
    boxed title style={
        colback=persianblue,
        colframe=persianblue
    },
    #1
}



%%%%%%%%%%%%%%%%%%%%%%%%%%%%%
%%%%% \begin{propriete} %%%%%
%%%%%%%%%%%%%%%%%%%%%%%%%%%%%

\newtcolorbox{propriete}[2][]{%
    enhanced , breakable , parbox=false ,
    arc=5pt,
    frame style={left color=persianblue, right color=softblue},
    colback=lightblue,
    fonttitle=\bfseries,
    title={\faBook~~ \underline{Propriété} : #2~~\faBook},
    attach boxed title to top left={xshift=5mm,yshift=-2mm},
    boxed title style={
        colback=persianblue,
        colframe=persianblue
    },
    #1
}



%%%%%%%%%%%%%%%%%%%%%%%%%%%%%%%%%%%%%%%%%%
%%%%% ----- \begin{definition} ----- %%%%%
%%%%%%%%%%%%%%%%%%%%%%%%%%%%%%%%%%%%%%%%%%

\newtcolorbox{definition}[2][]{%
    enhanced , breakable , parbox=false ,
    arc=5pt,
    frame style={left color=calgreen, right color=softgreen},
    colback=lightgreen,
    fonttitle=\bfseries,
    title={\faUniversity~~ \underline{Définition} : #2~~\faUniversity},
    attach boxed title to top left={xshift=5mm,yshift=-2mm},
    boxed title style={
        colback=calgreen,
        colframe=calgreen
    },
    #1
}



%%%%%%%%%%%%%%%%%%%%%%%%%%%%%%%%%%%%%%
%%%%% ----- \begin{astuce} ----- %%%%%
%%%%%%%%%%%%%%%%%%%%%%%%%%%%%%%%%%%%%%

\newtcolorbox{astuce}[1][]{%
    enhanced , parbox=false ,
    arc=10pt,
    frame hidden,
    colback=lightred,
    fonttitle=\bfseries,
    borderline={1mm}{0mm}{maroon,dotted},
    title={\faLightbulb~~\underline{Astuce}~~\faLightbulb},
    attach boxed title to top left={xshift=5mm,yshift=-2mm},
    boxed title style={
        colback=maroon,
        colframe=maroon
    },
    #1
}



%%%%%%%%%%%%%%%%%%%%%%%%%%%%%%%%%%%%%%%%
%%%%% ----- \begin{remarque} ----- %%%%%
%%%%%%%%%%%%%%%%%%%%%%%%%%%%%%%%%%%%%%%%

\newtcolorbox{remarque}[1][]{%
    enhanced , parbox=false ,
    arc=10pt,
    frame hidden,
    colback=lightred,
    fonttitle=\bfseries,
    borderline={1mm}{0mm}{maroon,dotted},
    title={\faExclamation~~\underline{Remarque}~~\faExclamation},
    attach boxed title to top left={xshift=5mm,yshift=-2mm},
    boxed title style={
        colback=maroon,
        colframe=maroon
    },
    #1
}



%%%%%%%%%%%%%%%%%%%%%%%%%%%%%%%%%%%%%%%
%%%%% ----- \begin{methode} ----- %%%%%
%%%%%%%%%%%%%%%%%%%%%%%%%%%%%%%%%%%%%%%

\newtcolorbox{methode}[2][]{%
    enhanced , parbox=false ,
    arc=10pt,
    frame hidden,
    colback=lightred,
    fonttitle=\bfseries,
    borderline={1mm}{0mm}{maroon,dotted},
    title={\faCogs~~\underline{Méthode} : #2~~\faCogs},
    attach boxed title to top left={xshift=5mm,yshift=-2mm},
    boxed title style={
        colback=maroon,
        colframe=maroon
    },
    #1
}



%%%%%%%%%%%%%%%%%%%%%%%%%%%%%%%%%%%%%%%
%%%%% ----- \begin{exemple} ----- %%%%%
%%%%%%%%%%%%%%%%%%%%%%%%%%%%%%%%%%%%%%%

\newcommand{\exempledecoration}{%
    \draw[color=brickorange , line width = 2pt , fill=white]
        (frame.south west) ++(0,1pt) -- ++(3cm,0) ++(0,1.5mm) rectangle ++ (3mm,-3mm) ;
}

\newtcolorbox{exemplebox}[1][]{%
	enhanced , breakable , parbox=false ,
    arc=0pt,
    frame hidden,
    borderline west={2pt}{0pt}{brickorange},
    interior style={
        top color=lightorange,
        bottom color=lightorange!15!white
    },
    fonttitle=\bfseries,
    overlay unbroken={\exempledecoration},
    overlay last={\exempledecoration},
    #1
}

\newenvironment{exemple}[1]{%
    \begin{exemplebox}
        \textbf{\color{brickorange}\underline{Exemple} : #1}\\[5pt]
}{%
    \end{exemplebox}
}




%%%%%%%%%%%%%%%%%%%%%%%%%%%%%%%%%%%%%%%%%%%%%
%%%%% ----- \begin{demonstration} ----- %%%%%
%%%%%%%%%%%%%%%%%%%%%%%%%%%%%%%%%%%%%%%%%%%%%

\newcommand{\demodecoration}{%
    \draw[color=borderpink , line width = 2pt , fill=white]
        (frame.south west) ++(0,1pt) -- ++(3cm,0) ++(0,1.5mm) rectangle ++ (3mm,-3mm) ;
}

\newtcolorbox{demobox}[1][]{%
	enhanced , breakable , parbox=false ,
    arc=0pt,
    frame hidden,
    borderline west={2pt}{0pt}{borderpink},
    interior style={
        top color=lightpink,
        bottom color=lightpink!15!white
    },
    fonttitle=\bfseries,
    overlay unbroken={\demodecoration},
    overlay last={\demodecoration},
    #1
}

\newenvironment{demonstration}{%
    \begin{demobox}
        \textbf{\color{borderpink}\underline{Démonstration}}\\[5pt]
}{%
    \end{demobox}
}



%%%%%%%%%%%%%%%%%%%%%%%%%%%%%%%%%
%%%%% ----- \exercice ----- %%%%%
%%%%%%%%%%%%%%%%%%%%%%%%%%%%%%%%%

\newtcolorbox[auto counter , number within=section , number freestyle={\noexpand\arabic{\tcbcounter}} ]{exercicebox}[2][]{% ----- Exercice box
enhanced,
arc=1pt,
colframe=black!80,
colback=black!10,
fonttitle=\bfseries,
hbox,
title={Exercice n° \thetcbcounter \, : #2 },
attach boxed title to top*,
boxed title style = {colback = black!80},
#1}


\newcommand\exercice[2]{% ----- Exercice command
    \begin{exercicebox}[]{#1}
        \begin{varwidth}{\linewidth-30pt} 
         #2
        \end{varwidth}
    \end{exercicebox}
}



%%%%%%%%%%%%%%%%%%%%%%%%%%%%%%%%
%%%%% ----- \calculs ----- %%%%%
%%%%%%%%%%%%%%%%%%%%%%%%%%%%%%%%

\newtcolorbox{calculbox}[1][]{% ----- Calculs box
enhanced,
arc=1pt,
colbacktitle=black!25!white,
coltitle=black,
colback=white,
colframe=black!25!white,
hbox,
title={Aides de calcul},
attach boxed title to top*,
#1}


\newcommand\calculs[1]{% ----- Calculs command
    \begin{calculbox}
        \restoreskip
        #1
    \end{calculbox}
}




%%%%%%%%%%%%%%%%%%%%%%%%%%%%%%%%%%%%%%
%%%%% ----- \begin{rappel} ----- %%%%%
%%%%%%%%%%%%%%%%%%%%%%%%%%%%%%%%%%%%%%

\newtcolorbox{rappel}[2][]{%
    enhanced , breakable , parbox=false ,
    arc=5pt,
    frame hidden,
    colback=lightgreen,
    fonttitle=\bfseries,
    borderline={2pt}{0mm}{calgreen , dashed},
    title={\faUniversity~~ \underline{Rappel} : #2~~\faUniversity},
    attach boxed title to top left={xshift=5mm,yshift=-2mm},
    boxed title style={
        colback=calgreen,
        colframe=calgreen
    },
    #1
}





%%%%%%%%%%%%%%%%%%%%%%%%
%%%%% Mise en page %%%%%
%%%%%%%%%%%%%%%%%%%%%%%%

%%%%% Établissement scolaire %%%%%
\def\etablissement{Lycée Polyvalent Irène et Frédéric Joliot Curie}



%%%%% Flèche SSI longueur adaptative %%%%%

\makeatletter

\newcommand{\ssi}[1]{%
  \@ifnextchar[{\ssi@i{#1}}{\ssi@i{#1}[]}%
}

\def\ssi@i#1[#2]{%
  \@ifnextchar[{\ssi@ii{#1}{#2}}{\ssi@ii{#1}{#2}[]}%
}

\def\ssi@ii#1#2[#3]{%
  \mathrel{%
    \ext@arrow 0099{%
      \Longleftrightarrowfill@
    }%
    {\makebox[#1]{#2}}%
    {\makebox[#1]{#3}}%
  }%
}

\def\Longleftrightarrowfill@{%
  \arrowfill@\Leftarrow\Relbar\Rightarrow
}

\makeatother



%%%%% Adaptation du stretch des environnements 'cases' %%%%%

\makeatletter
\renewcommand*\env@cases[1][1.2]{%
  \let\@ifnextchar\new@ifnextchar
  \left\lbrace
  \def\arraystretch{#1}%
  \array{@{}l@{\quad}l@{}}%
}
\makeatother



%%%%% Ajustement hauteur de \ldots %%%%%
\let\oldldots\ldots
\renewcommand{\ldots}{%
	\parbox[b][15pt]{15pt}{\oldldots}
}


%%%%% Colonne centrée de taille ajustable %%%%%
\newcolumntype{P}[1]{>{\centering\arraybackslash}m{#1}<{\vspace{0pt}}} 		   	% ----- Macros de mise en page

\geometry{left=30pt, right=30pt, top=75pt, bottom=55pt} % ----- Marges

\begin{document}


\firstpageheader{Prénom : ........................................\\[3mm] Nom : ............................................. }{}{Classe : T\up{ale} STMG4 \\[3mm] Date : ...............}
\footer{}{\footnotesize \etablissement}{\thepage / \numpages}

\hrulefill\raisebox{-1.9pt}{\quad\decofourleft~\decosix~\decofourright \quad}\hrulefill

\begin{center}
	\textbf{\large Devoir surveillé \no 1 (55 min)}
\end{center}
\vspace{-5pt}
Vous êtes invité\textperiodcentered e\textperiodcentered s à faire figurer sur votre copie toute trace de recherche, même incomplète ou non fructueuse, que vous aurez développée. La qualité de votre rédaction, la clarté et la précision de vos raisonnements seront prises en compte dans l’appréciation de votre copie.

Les élèves disposant d'un P.A.P. n'ont pas à traiter les questions marquées du symbole $\bigstar$

\vspace{-8pt}\hrulefill\raisebox{-1.9pt}{\quad\decofourleft~\decosix~\decofourright \quad}\hrulefill
\vspace{5pt}


%%%%%%%%%%%%%%%%%%%%%%%%%%%%%%%%%%%%%%%%%%%%%%%%%%%%%%%%%%%%%%%%%%%%%%%%%%%%%%%%%%%%%%%%%%%%%%%%%%%%%%%%%%%%%
%%%%%%%%%%%%%%%%%%%%%%%%%%%%%%%%%%%%%%%%%%%%%%%%%%%%%%%%%%%%%%%%%%%%%%%%%%%%%%%%%%%%%%%%%%%%%%%%%%%%%%%%%%%%%
%%%%%%%%%%%%%%%%%%%%%%%%%%%%%%%%%%%%%%%%%%%%%%%%%%%%%%%%%%%%%%%%%%%%%%%%%%%%%%%%%%%%%%%%%%%%%%%%%%%%%%%%%%%%%

\section{Questions de cours}
\begin{enumerate}[itemsep = 0pt]
	\item $\bigstar$ Écrire la formule du crible.
	\item Donner la formule des probabilités conditionnelles.
	\item À quelle condition peut-on affirmer que deux évènements $A$ et $B$ sont indépendants ?
\end{enumerate}



\section{Dans une urne}
Dans une urne on place 50 boules numérotées de $1$ à $50$. On considère les évènements suivants :

$\bullet$ A : \og La boule tirée porte un numéro pair \fg \\
$\bullet$ B : \og La boule tirée porte un numéro qui est un multiple de 5\fg \\
$\bullet$ C : \og La boule tirée porte un numéro supérieur ou égal à $30$ \fg
\begin{enumerate}[itemsep = 0pt]
	\item Déterminer $\PP(A)$ , $\PP(B)$ et $\PP(C)$
	\item $\bigstar$ Décrire explicitement l'évènement $A \cap B$
	\item \begin{enumerate}
		\item En justifiant, déterminer la valeur de $\PP(B \cap C)$
		\item En déduire alors $\PP_C(B)$
		\item Peut-on affirmer que les évènements $B$ et $C$ sont indépendants ?
	\end{enumerate}
\end{enumerate}



\section{État grippal}
Lors d'une épidémie de grippe, on vaccine un tiers de la population. On constate qu'un malade sur dix est vacciné et que la grippe touche un quart de la population.

On choisit un individu au hasard dans la population et on considère les évènements suivants :

$\bullet$ V : \og L'individu est vacciné \fg \\
$\bullet$ M : \og L'individu est malade \fg 
\begin{enumerate}[align = left]
	\item Grâce aux informations de l'énoncé et sans calcul, déterminer $\PP(M)$ , $\PP(V)$ et $\PP_M(V)$
	\item Les évènements $M$ et $V$ sont-ils indépendants ?
	\item $\bigstar$ Sur votre copie, construire un arbre pondéré permettant de représenter la situation et y placer, si possible, les probabilités de la question précédente.
	\item Calculer $\PP(M \cap V)$ et interpréter ce résultat dans le contexte de l'exercice.
	\item Dans cette question seulement, on suppose que l'individu choisi est vacciné.
	\begin{enumerate}
		\item Exprimer la probabilité que cet individu soit malade à l'aide des évènements $M$ et $V$.
		\item Calculer alors la probabilité qu'un individu vacciné soit malgré tout malade.
	\end{enumerate}
	\item[\underline{Question bonus} :] Déterminer $\PP(\barre{M} \cap V)$ , puis $\PP_{\barre{M}}(V)$
\end{enumerate}



%%%%%%%%%%%%%%%%%%%%%%%%%%%%%%%%%%%%%%%%%%%%%%%%%
\newpage
\newgeometry{left=30pt,right=30pt, top=30pt, bottom=55pt}
%%%%%%%%%%%%%%%%%%%%%%%%%%%%%%%%%%%%%%%%%%%%%%%%%



\section{D'un tableau naît un arbre}
On considère le tableau croisé d'effectifs suivant :
\begin{center}
	\renewcommand{\arraystretch}{1.75}
    \begin{tabular}{| P{4cm} | P{3cm} | P{3cm} | P{3cm} |}\hline
		\cellcolor{gray} & $A$ & $\barre{A}$ & Total\\ \hline
	    $B$ & 80 &  &  \\ \hline
	    $\barre{B}$ &  &  & 140 \\ \hline
	    Total & 120 &  & 200 \\ \hline
    \end{tabular}
\end{center}
\begin{enumerate}
	\item Compléter le tableau ci-dessus.
	\item Déterminer la valeur de $\PP_A(B)$ et de $\PP_B(A)$ en explicitant vos calculs.
	\item Compléter les deux arbres ci-après :
	
		\hspace{1cm}
		\begin{minipage}{0.45\textwidth}
			\begin{tikzpicture}[
				grow=right, scale = 0.8 ,
				level 1/.style={level distance=4cm, sibling distance=3cm},
				level 2/.style={level distance=3cm, sibling distance=1.5cm}
				]
				\small
				
				\node {$\Omega$}
				child {node {$\barre{A}$} 
					child {node {$\barre{B}$}
						edge from parent node[below] {\ldots}
					}
					child {node {$B$}
						edge from parent node[above] {\ldots}
					} edge from parent node[below] {\ldots}
				}
				child {node {$A$} 
					child {node {$\barre{B}$}
						edge from parent node[below] {\ldots}
					}
					child {node {$B$}
						edge from parent node[above] {\ldots}
					} edge from parent node[above] {\ldots}
				};
			\end{tikzpicture}
		\end{minipage}
		\begin{minipage}{0.5\textwidth}
			\begin{tikzpicture}[
				grow=right, scale = 0.8 ,
				level 1/.style={level distance=4cm, sibling distance=3cm},
				level 2/.style={level distance=3cm, sibling distance=1.5cm}
				]
				\small
				
				\node {$\Omega$}
				child {node {$\barre{B}$} 
					child {node {$\barre{A}$}
						edge from parent node[below] {\ldots}
					}
					child {node {$A$}
						edge from parent node[above] {\ldots}
					} edge from parent node[below] {\ldots}
				}
				child {node {$B$} 
					child {node {$\barre{A}$}
						edge from parent node[below] {\ldots}
					}
					child {node {$A$}
						edge from parent node[above] {\ldots}
					} edge from parent node[above] {\ldots}
				};
			\end{tikzpicture}
		\end{minipage}
	\item $\bigstar$ Déterminer la valeur de $\PP(A \cap B)$ -	\textit{\small Point bonus si vous utilisez plusieurs méthodes}
\end{enumerate}



\section{Q.C.M. (À points négatifs)}
Pour chacune des questions suivantes, une et une seule des réponses proposées est exacte. Pour répondre, entourez sur la feuille la réponse choisie. Aucune justification n'est demandée. \\
Une réponse correcte rapporte des points. L'absence de réponse ne rapporte ni ne retire aucun point.
Une réponse incorrecte ou une réponse multiple retire des points. . 

\begin{enumerate}[label=\underline{Question \no\arabic*} :  , leftmargin=* , itemsep=1mm]
	\item $\bigstar$ Si les évènements $A$ et $B$ sont indépendants et que $\PP(A) = 0,5$  et $\PP(B) = 0,4$ , combien vaut  $\PP(A \cap B)$ ?
	
		\hspace{\linewidth-\textwidth}
		\renewcommand{\arraystretch}{1.5}
		\begin{tabular}{| P{0.22\textwidth} | P{0.22\textwidth} | P{0.22\textwidth} | P{0.22\textwidth} |} \hline
			A & B & C & D \\ \hline
			$0,1$ & $0,2$ & $0,3$  & $0,4$ \\ \hline
		\end{tabular}
		
	\item Sur la route de son travail, Pierrick rencontre deux feux tricolores. Il a déterminé que la probabilité de devoir d'arrêter au premier feu est de $\dfrac{1}{3}$ et que celle de devoir s'arrêter au second est de $\dfrac{2}{5}$. Il a également déterminer qu'il y a une probabilité de $\dfrac{1}{2}$ de devoir s'arrêter à au moins un des deux feux tricolores. Quelle est la probabilité de devoir s'arrêter aux deux feux ?
	
		\hspace{\linewidth-\textwidth}
		\renewcommand{\arraystretch}{2}
		\begin{tabular}{| P{0.22\textwidth} | P{0.22\textwidth} | P{0.22\textwidth} | P{0.22\textwidth} |} \hline
			A & B & C & D \\ \hline
			$\dfrac{1}{30}$ & $\dfrac{7}{30}$ & $\dfrac{11}{30}$ & $\dfrac{13}{30}$ \\ \hline
		\end{tabular}
		
\end{enumerate}



\end{document}